\documentclass[11pt]{article}

\usepackage[margin=1in]{geometry}
\usepackage{graphicx}
\usepackage{amsmath}
\usepackage{siunitx}
\usepackage{caption}
\usepackage{fancyhdr}
\usepackage{float}
\usepackage{hyperref}
\usepackage{enumitem}

\pagestyle{fancy}
\setlength{\headheight}{14pt}
\fancyhf{}
\lhead{Project \#5: Circuit Design \#2}
\rhead{\thepage}
\renewcommand{\headrulewidth}{0.4pt}

\title{\vspace{-2cm}Linear Circuits 2 \\ Project \#5 Circuit Design \#2 Lab Report}
\author{Group Members: \\ Kenneth Gallmon \\ Juvon Mondesir}
\date{}

\begin{document}

\maketitle
\thispagestyle{fancy}

\section*{Objective}

The goal of the project is to design a circuit that meets certain design
requirements.

\begin{figure}[H]
    \centering
    \includegraphics[width=0.9\textwidth]{figures/fig01_circuit_model.png}
    \caption{Circuit model: function generator ($V_{IS}$, $R_S$) driving the
    circuit design under a variable load $R_L$}
    \label{fig:circuit_model}
\end{figure}

Basing our design off this circuit model, we are tasked to meet these
requirements:

\begin{figure}[H]
    \centering
    \includegraphics[width=0.6\textwidth]{figures/fig02_design_specs.png}
    \caption{Design specifications: resistance range, frequency-dependent
    gain requirements, and notations}
\end{figure}

So not only must we meet all these specifications, but the design must also
work for the entire range of load resistances.

\section*{Equipment}

\begin{itemize}[nosep]
    \item Lab Equipment
    \item Circuit Simulator
    \item Hantek 3-in-1 Digital Equipment
    \item Powered Breadboard
\end{itemize}

\section*{Simulations / Hand Calculations}

For all the simulations 10\,V is used in a sinusoidal voltage source over
varying frequencies. In addition to these, 50\,$\Omega$ is added in series
to replicate the function generator.

\subsection*{Initial Filter Behavior Simulation}

\begin{figure}[H]
    \centering
    \includegraphics[width=0.8\textwidth]{figures/fig03_initial_sim_rl1k.png}
    \caption{Frequency response and schematic with $R_2 = 1$\,k$\Omega$}
\end{figure}

\begin{figure}[H]
    \centering
    \includegraphics[width=0.55\textwidth]{figures/fig04_initial_sim_rl30k.png}
    \caption{Frequency response and schematic with $R_2 = 30$\,k$\Omega$}
\end{figure}

\textbf{Simulation 1)} $R_L = 30$\,k$\Omega$, $f = 1$\,kHz, $C = 0.047\,\mu$F

Using these values in the circuit equates to $|V_O/V_{IS}|$ being $0.275$.

\begin{figure}[H]
    \centering
    \includegraphics[width=0.4\textwidth]{figures/fig05_sim1_schematic.png}
    \caption{Simulation 1 schematic}
\end{figure}

\begin{figure}[H]
    \centering
    \includegraphics[width=0.6\textwidth]{figures/fig06_sim1_waveform.png}
    \caption{Simulation 1 output waveform}
\end{figure}

Above, is the simulated circuit and results of the gain, which align with
the specification of $|V_O| \leq 0.7|V_{IS}|$.

\textbf{Simulation 2)} $R_L = 1$\,k$\Omega$, $f = 10$\,kHz, $C = 0.047\,\mu$F

With these values the gain $|V_O/V_{IS}|$ of the circuit equates to $0.774$.

\begin{figure}[H]
    \centering
    \includegraphics[width=0.4\textwidth]{figures/fig07_sim2_schematic.png}
    \caption{Simulation 2 schematic}
\end{figure}

\begin{figure}[H]
    \centering
    \includegraphics[width=0.6\textwidth]{figures/fig08_sim2_waveform.png}
    \caption{Simulation 2 output waveform over a 10\,kHz span}
\end{figure}

\begin{figure}[H]
    \centering
    \includegraphics[width=0.9\textwidth]{figures/fig09_sim2_bode.png}
    \caption{Simulation 2 frequency response}
\end{figure}

Once again, the pictures above demonstrate the simulated circuit and the
waveform of $V_O$ over the span 10\,kHz. The gain of this circuit
$|V_O/V_{IS}|$ equates to $0.774$, which aligns with the specification of
$|V_O| \geq 0.7|V_{IS}|$.

\textbf{Simulation 3)} $R_L = 1$\,k$\Omega$, $f = 20$\,kHz, $C = 0.047\,\mu$F

\begin{figure}[H]
    \centering
    \includegraphics[width=0.4\textwidth]{figures/fig10_sim3_schematic.png}
    \caption{Simulation 3 schematic}
\end{figure}

\begin{figure}[H]
    \centering
    \includegraphics[width=0.6\textwidth]{figures/fig11_sim3_waveform.png}
    \caption{Simulation 3 output waveform over a 20\,kHz span}
\end{figure}

This is the 3rd simulated circuit and its waveform over the span of
20\,kHz. The gain of this circuit equates to $0.875$, which follows the
specification of this portion being $|V_O| \geq 0.85|V_{IS}|$.

\subsection*{Additional Hand Calculations}

\begin{figure}[H]
    \centering
    \includegraphics[width=0.5\textwidth]{figures/fig12_handcalc1.png}
    \caption{Hand calculations: equivalent resistance $R_{eq} = R \parallel
    R_L$ for the minimum and maximum load resistance}
\end{figure}

\begin{figure}[H]
    \centering
    \includegraphics[width=0.45\textwidth]{figures/fig13_handcalc2.png}
    \caption{Hand calculations verifying the gain requirements for
    $1\,\text{Hz} \leq f \leq 1\,\text{kHz}$ and $10\,\text{kHz} \leq f 
    20\,\text{kHz}$}
\end{figure}

\begin{figure}[H]
    \centering
    \includegraphics[width=0.5\textwidth]{figures/fig14_handcalc3.png}
    \caption{Hand calculations verifying the gain requirement for $f >
    20$\,kHz}
\end{figure}

\section*{Experiment}

To design the circuit, we knew it had to be a high pass filter, so we knew
to use a capacitor and resistor. As far as the values, after calculation
provided above, we concluded that a capacitor value of $0.047\,\mu$F and
internal resistor of 1\,k$\Omega$ would be sufficient to satisfy all
requirements, with minimal difference on varying load resistances as well.

\subsection*{Bode Plot Behavior}

\begin{figure}[H]
    \centering
    \includegraphics[width=0.95\textwidth]{figures/fig15_bode_plot_behavior.png}
    \caption{Measured Bode plot behavior of the filter}
\end{figure}

\subsection*{Scenario 1 Waveform: $R_L = 30$\,k$\Omega$, $f = 1$\,kHz}

\begin{figure}[H]
    \centering
    \includegraphics[width=0.85\textwidth]{figures/fig16_scenario1_waveform.png}
    \caption{Scenario 1 measured waveform (C1: input, C2: output)}
\end{figure}

$|V_O/V_{IS}|$ with scenario 1's values were $0.275$. In the picture above
from the oscilloscope, C1 is the voltage input signal and C2 is the voltage
output signal. So, comparing both signals it accurately reflects the
initial requirement of $|V_O| \leq 0.7|V_{IS}|$.

\subsection*{Scenario 2 Waveform: $R_L = 1$\,k$\Omega$, $f = 10$\,kHz}

\begin{figure}[H]
    \centering
    \includegraphics[width=0.85\textwidth]{figures/fig17_scenario2_waveform.png}
    \caption{Scenario 2 measured waveform (C1: input, C2: output)}
\end{figure}

In scenario 2, $|V_O/V_{IS}| = 0.774$. This image from the oscilloscope has
the same concept as the previous one, with Channel 1 being input and
Channel 2 being output. After comparison and viewing the $V_{pp}$ values you
can see these waveforms satisfy the requirement of $|V_O| \geq 0.7|V_{IS}|$.

\subsection*{Scenario 3 Waveform: $R_L = 1$\,k$\Omega$, $f = 20$\,kHz}

\begin{figure}[H]
    \centering
    \includegraphics[width=0.85\textwidth]{figures/fig18_scenario3_waveform.png}
    \caption{Scenario 3 measured waveform (C1: input, C2: output)}
\end{figure}

Lastly, in scenario 3, the $|V_O/V_{IS}| = 0.875$. From observing the
oscilloscope, it shows the correct ratio of $|V_O| \geq 0.85|V_{IS}|$.

\section*{Results}

Our output voltage matched all the constraints of the project, having
values close to $0.875\,V_{IS}$ for higher frequencies and close to
$0.275\,V_{IS}$ for lower frequencies. All in all, our filter circuit
performed well, and the results are as expected.

\section*{Conclusion}

Our circuit was a success, and we were able to use the properties of
capacitors under high and low frequencies to make a filter under specific
design constraints. We also tested the impact of the loading effect of the
resistor in tandem with the internal resistance of the function generator
and oscilloscope on filter response.

\end{document}